\documentclass[10pt,a4paper]{amsart}
\usepackage{amsmath,amssymb,amsthm,graphicx}
\usepackage{verbatim}
\newenvironment{reference}{\comment}{\endcomment}
%\newenvironment{reference}{}{}
\newenvironment{slogan}{\comment}{\endcomment}
\newenvironment{history}{\comment}{\endcomment}
% For commutative diagrams we use Xy-pic
\usepackage[all]{xy}

% We use 2cell for 2-commutative diagrams.
\xyoption{2cell}
% Theorem environments.
%
\theoremstyle{plain}
\newtheorem{theorem}[subsection]{Theorem}
\newtheorem{proposition}[subsection]{Proposition}
\newtheorem{lemma}[subsection]{Lemma}

\theoremstyle{definition}
\newtheorem{definition}[subsection]{Definition}
\newtheorem{example}[subsection]{Example}
\newtheorem{exercise}[subsection]{Exercise}
\newtheorem{situation}[subsection]{Situation}

\theoremstyle{remark}
\newtheorem{remark}[subsection]{Remark}
\newtheorem{remarks}[subsection]{Remarks}

\numberwithin{equation}{subsection}

% Macros
%
\def\lim{\mathop{\mathrm{lim}}\nolimits}
\def\colim{\mathop{\mathrm{colim}}\nolimits}
\def\Spec{\mathop{\mathrm{Spec}}}
\def\Hom{\mathop{\mathrm{Hom}}\nolimits}
\def\Ext{\mathop{\mathrm{Ext}}\nolimits}
\def\SheafHom{\mathop{\mathcal{H}\!\mathit{om}}\nolimits}
\def\SheafExt{\mathop{\mathcal{E}\!\mathit{xt}}\nolimits}
\def\Sch{\mathit{Sch}}
\def\Mor{\mathop{\mathrm{Mor}}\nolimits}
\def\Ob{\mathop{\mathrm{Ob}}\nolimits}
\def\Sh{\mathop{\mathit{Sh}}\nolimits}
\def\NL{\mathop{N\!L}\nolimits}
\def\CH{\mathop{\mathrm{CH}}\nolimits}
\def\proetale{{pro\text{-}\acute{e}tale}}
\def\etale{{\acute{e}tale}}
\def\QCoh{\mathit{QCoh}}
\def\Ker{\mathop{\mathrm{Ker}}}
\def\Im{\mathop{\mathrm{Im}}}
\def\Coker{\mathop{\mathrm{Coker}}}
\def\Coim{\mathop{\mathrm{Coim}}}

% Boxtimes
%
\DeclareMathSymbol{\boxtimes}{\mathbin}{AMSa}{"02}

%
% Macros for moduli stacks/spaces
%
\def\QCohstack{\mathcal{QC}\!\mathit{oh}}
\def\Cohstack{\mathcal{C}\!\mathit{oh}}
\def\Spacesstack{\mathcal{S}\!\mathit{paces}}
\def\Quotfunctor{\mathrm{Quot}}
\def\Hilbfunctor{\mathrm{Hilb}}
\def\Curvesstack{\mathcal{C}\!\mathit{urves}}
\def\Polarizedstack{\mathcal{P}\!\mathit{olarized}}
\def\Complexesstack{\mathcal{C}\!\mathit{omplexes}}
% \Pic is the operator that assigns to X its picard group, usage \Pic(X)
% \Picardstack_{X/B} denotes the Picard stack of X over B
% \Picardfunctor_{X/B} denotes the Picard functor of X over B
\def\Pic{\mathop{\mathrm{Pic}}\nolimits}
\def\Picardstack{\mathcal{P}\!\mathit{ic}}
\def\Picardfunctor{\mathrm{Pic}}
\def\Deformationcategory{\mathcal{D}\!\mathit{ef}}
\usepackage{cleveref}
\crefname{originallemma}{original lemma}{original lemmas}
\Crefname{originallemma}{Original lemma}{Original lemmas}
\makeatletter
\newlabel{lemma-compose-after-separated}{{lemma-compose-after-separated}{source}{}{sourcepointer-1}{}}
\newlabel{lemma-compose-after-separated@cref}{{[originallemma][0][]lemma-compose-after-separated}{[1][0][]source}{}{}{}}
\makeatother
\crefname{originalreference}{original reference}{original references}
\Crefname{originalreference}{Original reference}{Original references}
\makeatletter
\newlabel{spaces-definition-separated}{{spaces-definition-separated}{source}{}{sourcepointer-2}{}}
\newlabel{spaces-definition-separated@cref}{{[originalreference][0][]spaces-definition-separated}{[1][0][]source}{}{}{}}
\makeatother
\crefname{originalreference}{original reference}{original references}
\Crefname{originalreference}{Original reference}{Original references}
\makeatletter
\newlabel{spaces-morphisms-lemma-composition-quasi-finite}{{spaces-morphisms-lemma-composition-quasi-finite}{source}{}{sourcepointer-3}{}}
\newlabel{spaces-morphisms-lemma-composition-quasi-finite@cref}{{[originalreference][0][]spaces-morphisms-lemma-composition-quasi-finite}{[1][0][]source}{}{}{}}
\makeatother
\crefname{originalreference}{original reference}{original references}
\Crefname{originalreference}{Original reference}{Original references}
\makeatletter
\newlabel{spaces-morphisms-lemma-composition-unramified}{{spaces-morphisms-lemma-composition-unramified}{source}{}{sourcepointer-4}{}}
\newlabel{spaces-morphisms-lemma-composition-unramified@cref}{{[originalreference][0][]spaces-morphisms-lemma-composition-unramified}{[1][0][]source}{}{}{}}
\makeatother
\crefname{originalreference}{original reference}{original references}
\Crefname{originalreference}{Original reference}{Original references}
\makeatletter
\newlabel{spaces-morphisms-lemma-permanence-quasi-finite}{{spaces-morphisms-lemma-permanence-quasi-finite}{source}{}{sourcepointer-5}{}}
\newlabel{spaces-morphisms-lemma-permanence-quasi-finite@cref}{{[originalreference][0][]spaces-morphisms-lemma-permanence-quasi-finite}{[1][0][]source}{}{}{}}
\makeatother
\crefname{originalreference}{original reference}{original references}
\Crefname{originalreference}{Original reference}{Original references}
\makeatletter
\newlabel{spaces-morphisms-lemma-permanence-unramified}{{spaces-morphisms-lemma-permanence-unramified}{source}{}{sourcepointer-6}{}}
\newlabel{spaces-morphisms-lemma-permanence-unramified@cref}{{[originalreference][0][]spaces-morphisms-lemma-permanence-unramified}{[1][0][]source}{}{}{}}
\makeatother
\title{A stack separated over the base has separated diagonal and consequently separated structure morphism}
\author{The Stacks Project Authors}
\date{stacks-project:stacks-morphisms.tex:local-snapshot-20261012; GNU FDL 1.2 or later}
\begin{document}
\maketitle
\noindent\textit{Editorial note.} All delivered mathematical text is the author's own, copied from the pinned Stacks Project chapter source sources/stacks-algebraic-stacks-morphisms/pinned-original-source.tex; no statement or proof step was rewritten, reordered or removed.\par\medskip
\noindent\textit{Editorial note.} Preamble provenance: the Stacks Project keeps one single shared preamble for every chapter; this card takes it directly from the staged sources/stacks-algebraic-stacks-morphisms/preamble.tex as three verbatim spans (lines 8-12, 14-18 and 154-217), so no derived file was written and nothing was added or redefined. The wrapper class and the base packages are supplied by the assembler.\par\medskip
\noindent\textit{Editorial note.} Bibliography provenance: neither the counted statement nor its proof contains a citation (this lane rescans the whole span before writing the card), so the card carries no bibliography and no entry was reconstructed.\par\medskip
\noindent\textit{Editorial note.} Selection provenance: this card was opened under the file-based control contract (CONTROL.md), with the mechanical pre-card gate run on its statement and proof interval before the card was written.\par\medskip
\noindent\textit{Allowed standard prerequisites (editorial).}\begin{itemize}
\item The description of the diagonal of the structure morphism through the diagonal of the stack over the base. Used at the start of the counted proof (source lines 1088-1098). Identifies the two diagonals so that properties can be transferred.
\item The base change permanence and the characterizations of the separation properties through the diagonal. Invoked inside the counted proof (source lines 1099-1108). Moves the properties between the two diagonals.
\item The identity stack as a representable and separated object over the base. Used at the end of the counted proof (source lines 1109-1113). Provides the trivial case in the comparison.
\end{itemize}\medskip
\ifdefined\hypertarget\hypertarget{sourcepointer-1}{}\fi
\noindent\textit{Source reference.} Original lemma lemma-compose-after-separated: sources/stacks-algebraic-stacks-morphisms/pinned-original-source.tex, source lines 929--929. Navigation pointer only: the counted text refers to the Stacks item lemma-compose-after-separated, which is an existing labeled result of this same chapter and is not part of this card. The pointer carries the target's own label as its number and its definition line as its span; its content is not reproduced here.\par\medskip
\ifdefined\hypertarget\hypertarget{sourcepointer-2}{}\fi
\noindent\textit{Source reference.} Original reference spaces-definition-separated: sources/stacks-algebraic-stacks-morphisms/pinned-original-source.tex, source lines 1091--1091. Navigation pointer only: the counted text refers to the Stacks item spaces-definition-separated, which belongs to another chapter of the project and is not part of this card. The pointer carries the target's own label as its number because this chapter's staged snapshot does not contain the other chapter's tag table; its content is not reproduced here.\par\medskip
\ifdefined\hypertarget\hypertarget{sourcepointer-3}{}\fi
\noindent\textit{Source reference.} Original reference spaces-morphisms-lemma-composition-quasi-finite: sources/stacks-algebraic-stacks-morphisms/pinned-original-source.tex, source lines 881--881. Navigation pointer only: the counted text refers to the Stacks item spaces-morphisms-lemma-composition-quasi-finite, which belongs to another chapter of the project and is not part of this card. The pointer carries the target's own label as its number because this chapter's staged snapshot does not contain the other chapter's tag table; its content is not reproduced here.\par\medskip
\ifdefined\hypertarget\hypertarget{sourcepointer-4}{}\fi
\noindent\textit{Source reference.} Original reference spaces-morphisms-lemma-composition-unramified: sources/stacks-algebraic-stacks-morphisms/pinned-original-source.tex, source lines 880--880. Navigation pointer only: the counted text refers to the Stacks item spaces-morphisms-lemma-composition-unramified, which belongs to another chapter of the project and is not part of this card. The pointer carries the target's own label as its number because this chapter's staged snapshot does not contain the other chapter's tag table; its content is not reproduced here.\par\medskip
\ifdefined\hypertarget\hypertarget{sourcepointer-5}{}\fi
\noindent\textit{Source reference.} Original reference spaces-morphisms-lemma-permanence-quasi-finite: sources/stacks-algebraic-stacks-morphisms/pinned-original-source.tex, source lines 971--971. Navigation pointer only: the counted text refers to the Stacks item spaces-morphisms-lemma-permanence-quasi-finite, which belongs to another chapter of the project and is not part of this card. The pointer carries the target's own label as its number because this chapter's staged snapshot does not contain the other chapter's tag table; its content is not reproduced here.\par\medskip
\ifdefined\hypertarget\hypertarget{sourcepointer-6}{}\fi
\noindent\textit{Source reference.} Original reference spaces-morphisms-lemma-permanence-unramified: sources/stacks-algebraic-stacks-morphisms/pinned-original-source.tex, source lines 968--968. Navigation pointer only: the counted text refers to the Stacks item spaces-morphisms-lemma-permanence-unramified, which belongs to another chapter of the project and is not part of this card. The pointer carries the target's own label as its number because this chapter's staged snapshot does not contain the other chapter's tag table; its content is not reproduced here.\par\medskip
\begin{lemma}
\label{lemma-separated-implies-morphism-separated}
Let $\mathcal{X}$ be an algebraic stack over the base scheme $S$.
\begin{enumerate}
\item
$\mathcal{X}$ is DM $\Leftrightarrow$
$\mathcal{X}$ is DM over $S$.
\item
$\mathcal{X}$ is quasi-DM $\Leftrightarrow$
$\mathcal{X}$ is quasi-DM over $S$.
\item If $\mathcal{X}$ is separated, then
$\mathcal{X}$ is separated over $S$.
\item If $\mathcal{X}$ is quasi-separated, then
$\mathcal{X}$ is quasi-separated over $S$.
\end{enumerate}
Let $f : \mathcal{X} \to \mathcal{Y}$ be a morphism of algebraic stacks
over the base scheme $S$.
\begin{enumerate}
\item[(5)] If $\mathcal{X}$ is DM over $S$, then $f$ is DM.
\item[(6)] If $\mathcal{X}$ is quasi-DM over $S$, then $f$ is quasi-DM.
\item[(7)] If $\mathcal{X}$ is separated over $S$ and
$\Delta_{\mathcal{Y}/S}$ is separated, then $f$ is separated.
\item[(8)] If $\mathcal{X}$ is quasi-separated over $S$ and
$\Delta_{\mathcal{Y}/S}$ is quasi-separated, then $f$ is quasi-separated.
\end{enumerate}
\end{lemma}

\begin{proof}
Parts (5), (6), (7), and (8) follow immediately from
Lemma \ref{lemma-compose-after-separated}
and
Spaces, Definition \ref{spaces-definition-separated}.
To prove (3) and (4) think of $X$ and $Y$ as algebraic stacks over
$\Spec(\mathbf{Z})$ and apply
Lemma \ref{lemma-compose-after-separated}.
Similarly, to prove (1) and (2), think of $\mathcal{X}$ as an algebraic
stack over $\Spec(\mathbf{Z})$ consider the
morphisms
$$
\mathcal{X} \longrightarrow
\mathcal{X} \times_S \mathcal{X} \longrightarrow
\mathcal{X} \times_{\Spec(\mathbf{Z})} \mathcal{X}
$$
Both arrows are representable by algebraic spaces.
The second arrow is unramified and locally quasi-finite as the base change of
the immersion $\Delta_{S/\mathbf{Z}}$. Hence the composition is
unramified (resp.\ locally quasi-finite) if and only if the first arrow
is unramified (resp.\ locally quasi-finite), see
Morphisms of Spaces,
Lemmas \ref{spaces-morphisms-lemma-composition-unramified} and
\ref{spaces-morphisms-lemma-permanence-unramified}
(resp.\ Morphisms of Spaces,
Lemmas \ref{spaces-morphisms-lemma-composition-quasi-finite} and
\ref{spaces-morphisms-lemma-permanence-quasi-finite}).
\end{proof}

\end{document}
